\documentclass[12pt,a4paper]{report}

\usepackage{fontspec}
\usepackage{polyglossia}
\setdefaultlanguage{romanian}
\usepackage[margin=2.25cm]{geometry}
\usepackage{setspace}
\usepackage{amsmath}
\usepackage{amssymb}
\usepackage{enumitem}
\usepackage{longtable}
\usepackage{array}
\usepackage{booktabs}
\usepackage{xcolor}
\usepackage{listings}
\usepackage[
    colorlinks,
    linkcolor={black},
    citecolor={black},
    urlcolor={blue}
]{hyperref}

\onehalfspacing
\setlength{\emergencystretch}{3em}
\setlength{\parskip}{0.15em}
\setlist{nosep,leftmargin=*}

\definecolor{codebackground}{rgb}{0.96,0.96,0.96}
\definecolor{codecomment}{rgb}{0.23,0.45,0.30}
\definecolor{codekeyword}{rgb}{0.10,0.20,0.55}
\definecolor{codestring}{rgb}{0.55,0.15,0.15}

\lstdefinestyle{codestyle}{
    backgroundcolor=\color{codebackground},
    basicstyle=\ttfamily\scriptsize,
    breaklines=true,
    frame=single,
    columns=fullflexible,
    keepspaces=true,
    showstringspaces=false,
    numbers=left,
    numberstyle=\tiny\color{gray},
    keywordstyle=\color{codekeyword}\bfseries,
    commentstyle=\color{codecomment},
    stringstyle=\color{codestring},
    tabsize=4
}
\lstset{style=codestyle}

\newcommand{\codechunk}[2]{%
\subsection*{Liniile #1--#2}
\lstinputlisting[language=Python,firstline=#1,lastline=#2,firstnumber=#1]{../src/dsp/spectral_subtraction.py}
}

\title{Explicarea fisierului \texttt{spectral\_subtraction.py}}
\author{}
\date{}

\begin{document}

\maketitle
\tableofcontents
\clearpage

\chapter{Rolul fisierului \texttt{spectral\_subtraction.py}}

Fisierul \path{src/dsp/spectral_subtraction.py} implementeaza metoda clasica
\textit{Spectral Subtraction}. Aceasta este una dintre metodele de referinta ale
proiectului, folosita pentru comparatia cu modelul neuronal MaskNet. Spre deosebire
de MaskNet, metoda nu are antrenare si nu invata din date, ci estimeaza un profil
de zgomot direct din semnalul zgomotos si il scade din spectru.

Modelul matematic este:

\[
y[n] = s[n] + v[n],
\]

unde \(y[n]\) este semnalul zgomotos, \(s[n]\) este vorbirea curata, iar \(v[n]\)
este zgomotul. In domeniul STFT:

\[
Y[k,m] = S[k,m] + V[k,m].
\]

Metoda estimeaza magnitudinea zgomotului \(\hat{N}[k]\), apoi calculeaza:

\[
|\hat{S}[k,m]| =
\max\left(|Y[k,m]| - \alpha\hat{N}[k], \beta\hat{N}[k]\right).
\]

\(\alpha\) controleaza cat de agresiv se scade zgomotul, iar \(\beta\) este
podeaua spectrala care previne valori negative sau reducere excesiva.

\chapter{Importuri}

\codechunk{1}{4}

\begin{description}
    \item[Linia 1.] Importa NumPy, folosit pentru calcule pe spectrograme complexe,
    energii de cadre, sortare si reconstructie magnitudine-faza.

    \item[Linia 2.] Importa PyTorch, folosit pentru conversia waveform-ului in tensor
    si pentru STFT/ISTFT prin utilitarele PyTorch ale proiectului.

    \item[Linia 3.] Importa \texttt{stft\_utils}, modulul care contine functiile
    comune pentru STFT, ISTFT, separare magnitudine/faza si reconstructie.

    \item[Linia 4.] Importa \texttt{config}, de unde sunt preluati parametrii STFT:
    \texttt{N\_FFT}, \texttt{HOP\_LEN} si \texttt{FRAME\_LEN}.
\end{description}

\chapter{Estimarea magnitudinii zgomotului}

\codechunk{6}{13}

\begin{description}
    \item[Linia 6.] Defineste functia \texttt{estimate\_noise\_mag}. Intrarea este
    spectrograma complexa a semnalului zgomotos, iar iesirea este o estimare a
    magnitudinii zgomotului pe fiecare bin de frecventa.

    \item[Linia 7.] Docstring-ul precizeaza ca zgomotul este estimat din cadrele
    cu energie minima, fara etichete VAD externe.

    \item[Linia 8.] Se separa magnitudinea si faza spectrogramei zgomotoase. Pentru
    estimarea zgomotului este folosita doar magnitudinea.

    \item[Linia 9.] \texttt{num\_frames} reprezinta numarul de cadre temporale din
    spectrograma.

    \item[Linia 10.] Energia fiecarui cadru este calculata ca media patratelor
    magnitudinilor:
    \[
    E[m] = \frac{1}{K}\sum_k |Y[k,m]|^2.
    \]

    \item[Linia 11.] Cadrele sunt sortate crescator dupa energie, iar functia alege
    primele \texttt{num\_noise\_frames}. Ipoteza este ca aceste cadre contin mai
    putina vorbire si mai mult zgomot de fundal.

    \item[Linia 12.] Magnitudinea zgomotului este media magnitudinilor din cadrele
    alese:
    \[
    \hat{N}[k] = \frac{1}{M}\sum_{m \in \mathcal{I}} |Y[k,m]|.
    \]
    \texttt{keepdims=True} pastreaza forma \((F,1)\), utila la broadcasting.

    \item[Linia 13.] Intoarce profilul estimat al zgomotului.
\end{description}

\chapter{Aplicarea metodei Spectral Subtraction}

\codechunk{15}{26}

\begin{description}
    \item[Linia 15.] Defineste functia \texttt{spectral\_subtraction}. Ea primeste
    spectrograma zgomotoasa, profilul de zgomot, factorul \(\alpha\) si podeaua
    spectrala \(\beta\).

    \item[Liniile 16--18.] Docstring-ul precizeaza ca metoda foloseste un profil
    de zgomot stationar estimat intern.

    \item[Linia 19.] Se separa magnitudinea si faza. Spectral Subtraction modifica
    magnitudinea, dar pastreaza faza semnalului zgomotos.

    \item[Liniile 20--21.] Daca zgomotul a fost estimat ca vector \((F,1)\), acesta
    este repetat pe toate cadrele pentru a avea aceeasi forma ca magnitudinea
    semnalului zgomotos.

    \item[Linia 22.] Se calculeaza magnitudinea dupa scaderea zgomotului:
    \[
    |Y[k,m]| - \alpha\hat{N}[k].
    \]

    \item[Linia 23.] Se calculeaza podeaua spectrala:
    \[
    \beta\hat{N}[k].
    \]
    Aceasta impiedica magnitudinea sa coboare sub o valoare minima proportionala
    cu zgomotul estimat.

    \item[Linia 24.] Magnitudinea imbunatatita este maximul dintre magnitudinea
    scazuta si podeaua spectrala:
    \[
    |\hat{S}[k,m]| =
    \max(|Y[k,m]|-\alpha\hat{N}[k], \beta\hat{N}[k]).
    \]

    \item[Linia 25.] Se reconstruieste spectrograma complexa folosind magnitudinea
    imbunatatita si faza initiala.

    \item[Linia 26.] Intoarce spectrograma complexa procesata.
\end{description}

\chapter{Functia completa pentru waveform}

\codechunk{28}{45}

\begin{description}
    \item[Linia 28.] \texttt{enhance\_waveform} este functia de nivel inalt care
    primeste un semnal temporal si intoarce semnalul imbunatatit.

    \item[Liniile 29--31.] Daca semnalul are mai multe canale, este convertit la
    mono prin mediere. Apoi este convertit la \texttt{float32}.

    \item[Liniile 33--40.] Semnalul temporal este transformat in STFT prin
    \texttt{stft\_utils.compute\_stft}. Sunt folositi parametrii centrali din
    \texttt{config.py}, pentru a pastra compatibilitatea cu restul proiectului.

    \item[Liniile 42--43.] Magnitudinea si faza returnate de STFT sunt convertite
    inapoi intr-o spectrograma complexa:
    \[
    Y[k,m] = |Y[k,m]|e^{j\angle Y[k,m]}.
    \]

    \item[Linia 44.] Se estimeaza profilul de zgomot din cadrele cu energie mica.

    \item[Linia 45.] Se aplica efectiv metoda Spectral Subtraction asupra spectrogramei.
\end{description}

\codechunk{47}{64}

\begin{description}
    \item[Liniile 47--49.] Din spectrograma imbunatatita sunt extrase magnitudinea
    si faza, pentru a putea folosi functia comuna \texttt{inverse\_stft}.

    \item[Liniile 51--53.] Magnitudinea si faza sunt convertite din NumPy in tensori
    PyTorch \texttt{float}.

    \item[Liniile 54--61.] Se reconstruieste waveform-ul prin ISTFT, folosind aceiasi
    parametri STFT si lungimea originala a semnalului. Parametrul \texttt{length}
    este important pentru ca iesirea sa aiba aceeasi lungime ca intrarea.

    \item[Liniile 63--64.] Tensorul rezultat este convertit inapoi in NumPy si
    returnat.
\end{description}

\chapter{Avantaje si limitari}

Spectral Subtraction este utila in proiect deoarece ofera un baseline usor de
explicat si rapid de rulat. Nu necesita training si poate functiona chiar daca
nu exista checkpoint MaskNet. Este folosita in demo-ul offline, in evaluarea
comparativa si in demo-ul real-time ca alternativa mai usoara decat modelul
neuronal.

Avantajele metodei sunt:

\begin{itemize}
    \item simplitate matematica;
    \item cost computational redus;
    \item lipsa dependentei de date de antrenare;
    \item interpretabilitate ridicata.
\end{itemize}

Limitarea principala este estimarea zgomotului. Daca primele cadre cu energie
mica nu reprezinta bine zgomotul real, scaderea poate fi prea slaba sau prea
agresiva. In plus, metoda poate produce artefacte muzicale, mai ales cand zgomotul
este variabil sau cand \(\alpha\) este prea mare.

\chapter{Concluzie}

\texttt{spectral\_subtraction.py} implementeaza un baseline clasic pentru
imbunatatirea vorbirii. Fluxul este clar: waveform \(\rightarrow\) STFT
\(\rightarrow\) estimare zgomot \(\rightarrow\) scadere spectrala \(\rightarrow\)
ISTFT. Aceasta metoda este importanta in proiect deoarece ofera un punct de
comparatie pentru MaskNet si pentru Wiener Filter.

\end{document}
